% FortySecondsCV LaTeX template
% Copyright © 2019 René Wirnata <rene.wirnata@pandascience.net>
% Licensed under the 3-Clause BSD License. See LICENSE file for details.
%
% Attributions
% ------------
% * fortysecondscv is based on the twentysecondcv class by Carmine Spagnuolo 
%   (cspagnuolo@unisa.it), released under the MIT license and available under
%   https://github.com/spagnuolocarmine/TwentySecondsCurriculumVitae-LaTex
% * further attributions are indicated immediately before corresponding code

%-------------------------------------------------------------------------------
%                             ADDITIONAL PACKAGES
%-------------------------------------------------------------------------------
\documentclass[
  a4paper, 
%   showframes,
   maincolor=cvblue,
   sectioncolor=cvblue,
%  subsectioncolor=orange
%   sidebarwidth=0.4\paperwidth,
%   topbottommargin=0.03\paperheight,
%   leftrightmargin=20pt
]{fortysecondscv}

% improve word spacing and hyphenation
\usepackage{microtype}
\usepackage{ragged2e}
\usepackage{rotating}

% take care of proper font encoding
\ifxetex
	\usepackage{fontspec}
	\defaultfontfeatures{Ligatures=TeX}
% \newfontfamily\headingfont[Path = fonts/]{segoeuib.ttf} % local font
\else
	\usepackage[utf8]{inputenc}
	\usepackage[T1]{fontenc}
% \usepackage[sfdefault]{noto} % use noto google font
\fi

% enable mathematical syntax for some symbols like \varnothing
\usepackage{amssymb}

% bubble diagram configuration
\usepackage{smartdiagram}
\smartdiagramset{
  % defaut font size is \large, so adjust to harmonize with sidebar layout
  bubble center node font = \footnotesize,
  bubble node font = \footnotesize,
  % default: 4cm/2.5cm; make minimum diameter relative to sidebar size
  bubble center node size = 0.3\sidebartextwidth,
  bubble node size = 0.25\sidebartextwidth,
  distance center/other bubbles = 1.5em,
  % set center bubble color
  bubble center node color = maincolor!70,
  % define the list of colors usable in the diagram
  set color list = {maincolor!10, maincolor!40,
  maincolor!20, maincolor!60, maincolor!35},
  % sets the opacity at which the bubbles are shown
  bubble fill opacity = 0.8,
}


%-------------------------------------------------------------------------------
%                            PERSONAL INFORMATION
%-------------------------------------------------------------------------------
% profile picture
\cvprofilepic{pics/AmirMahdiAmani}
% your name
\cvname{\LARGE Amir Mahdi Amani}
% job title/career
\cvjobtitle{AI Solution Engineer\\[0.2em]AI, Vision \& Automation}

% information row
\cvbirthday{\textbf{Location}: Cologne, Germany}
% short address/location, use \newline if more than 1 line is required
\cvaddress{\mbox{Bernhardstr.~119, 50968~Cologne}}
% phone number
\cvphone{+49-152-11555803}
% email address
\cvmail{amirmahdi.amani7th@gmail.com}
% professional profile
\cvsite{\href{https://www.linkedin.com/in/amirmahdiamani}{\textbf{LinkedIn}: amirmahdiamani}}
% GitHub profile (the original class macro is ORCID-oriented, so the display is reset below)
\cvkey{GitHub}{am-amani}
\renewcommand{\cvkey}{\href{https://github.com/am-amani}{\textbf{GitHub}: am-amani}}


%-------------------------------------------------------------------------------
%                              SIDEBAR 1st PAGE
%-------------------------------------------------------------------------------
% add more profile sections to sidebar on first page
\addtofrontsidebar{
	% include gosquare national flags from https://github.com/gosquared/flags;
	% naming according to ISO 3166-1 alpha-2 country codes
	\graphicspath{{pics/flags/}}

	\profilesection{Focus Areas}
		\skill{\faGraduationCap}{Multi-Sensor Navigation}
		\skill{\faGraduationCap}{Deep Reinforcement Learning}
		\skill{\faGraduationCap}{Computer Vision}
		\skill{\faGraduationCap}{Automation}
		\skill{\faGraduationCap}{Geospatial Data Analysis}

	\profilesection{Skills}
	\chartlabel{Programming:}\\
		\pointskill{}{Python}{5}
		\pointskill{}{C++}{4}
		\pointskill{}{SQL}{3}
		\pointskill{}{MATLAB, R}{3}
		\pointskill{}{Bash, LaTeX}{3}
	\chartlabel{AI and Vision:}\\
		\pointskill{}{PyTorch, TensorFlow}{4}
		\pointskill{}{OpenCV, YOLOv5}{4}
		\pointskill{}{Image Processing}{4}
	\chartlabel{Robotics:}\\
		\pointskill{}{ROS 1/2, Gazebo}{5}
		\pointskill{}{SLAM, AMCL, MoveIt}{4}
	\chartlabel{Tools and Systems:}\\
		\pointskill{}{Linux, Docker, Git}{4}
		\pointskill{}{QGIS, MS Office}{4}
}


%-------------------------------------------------------------------------------
%                              SIDEBAR 2nd PAGE
%-------------------------------------------------------------------------------
\addtobacksidebar{
	\profilesection{Short Bio}
	\aboutme{AI and robotics engineer with an M.Sc. in Computer Engineering, specializing in Artificial Intelligence and Robotics. Experienced in Python, C++, deep learning, computer vision, autonomous navigation, geospatial analysis, technical testing, and documentation. Research work includes multi-sensor robot navigation using teacher-student policy distillation in ROS/Gazebo. Currently developing data-driven solutions for network-planning and automation workflows in Germany.}

	\profilesection{Metrics}
	\begin{figure}\centering
		\smartdiagram[bubble diagram]{\large
			\textcolor{black!90}{GRE}\\
			\textcolor{black!90}{\Large 333},
			\textcolor{black!90}{Papers}\\
			\textcolor{black!90}{1},
			\textcolor{black!90}{Awards}\\
			\textcolor{black!90}{2},
			\textcolor{black!90}{Projects}\\
			\textcolor{black!90}{6}
		}
	\end{figure}

	\profilesection{Profiles}
	\centering
	\href{https://www.linkedin.com/in/amirmahdiamani}{\includegraphics[width=0.2\sidebartextwidth]{in.png}}\quad
	\href{https://github.com/am-amani}{\includegraphics[width=0.2\sidebartextwidth]{github.png}}

	\profilesection{Languages}
	\barskill{}{\textbf{English} (Professional)}{85}
	\barskill{}{\textbf{German} (B1)}{45}

	\profilesection{Personal}
	\begin{memberships}
		\membership{}{Interests include cooking, karaoke, dancing, reading, gym training, open discussions, and films. As a member of the Computer Science Association, helped organize an annual ACM programming competition for students from different universities in South Khorasan and conducted monthly C++ and Python workshops for freshmen. Also organized weekly English discussion sessions, special events, and Persian karaoke activities. Later created an online page for Persian karaoke content and community engagement.}
	\end{memberships}
}

\addtothirdsidebar{
\textbf{}\\
\begin{turn}{-90}\chartlabel{Multi-Sensor Navigation}\chartlabel{Deep Reinforcement Learning}\chartlabel{Computer Vision}\chartlabel{Autonomous Robots}\chartlabel{ROS and Gazebo}\chartlabel{Geospatial Analysis}\chartlabel{Python Automation}\end{turn}
}

%-------------------------------------------------------------------------------
%                         TABLE ENTRIES RIGHT COLUMN
%-------------------------------------------------------------------------------
\begin{document}

\newpage
\restoregeometry
\sidebarwidth=0.35\paperwidth

\makefrontsidebar

\cvsection{Working Experience}
\begin{cvtable}
	\cvitem{Oct. 2024 -- present}{\color{cvsectioncolor}AI Solution Engineer -- Network Planning, Optimization \& Automation Specialist}{\textbf{Deutsche Netzplanung GmbH, Cologne, Germany}}{Contribute to the development and optimization of telecommunications infrastructure. Use QGIS to analyze surface areas and classify regions for efficient fiber-network deployment. Collaborate with senior engineers on technical solutions, analyze project data, and prepare detailed reports and stakeholder presentations. The role combines data analysis, software-supported planning, and technical communication.}

	\cvitem{Mar. 2024 -- Sep. 2024}{\color{cvsectioncolor}Fiber Network Planning and Optimization (Part-Time)}{\textbf{Deutsche Netzplanung GmbH, Cologne, Germany}}{Supported network-planning projects and the optimization of fiber infrastructure. Used QGIS for surface-area analysis and regional classification, reviewed technical data, assisted engineers with planning decisions, and supported implementation and project documentation.}

	\cvitem{Feb. 2024 -- Sep. 2024}{\color{cvsectioncolor}Robotic Research Assistant (HiWi)}{\textbf{Humanoid Robots Lab, University of Bonn}}{Analyzed and documented the Unitree GO1 architecture and operating protocols. Installed and configured ROS packages including SLAM and AMCL, tested different operating modes, prepared manuals for researchers, and upgraded TurtleBot2 systems for localization, navigation, and mapping to support ongoing laboratory work.}

	\cvitem{Sep. 2019 -- Nov. 2019}{\color{cvsectioncolor}Research Internship -- Design-to-Robotic-Production and Operation}{\textbf{TU Delft, The Netherlands}}{Developed ROS-based solutions for research projects using MoveIt and Kinect V2. Applied forward and inverse kinematics for serial manipulators and implemented RGB-D object analysis through calibration, filtering, segmentation, clustering, object detection, and pose estimation. Gained practical exposure to the KUKA KR 10 and C4 Compact Controller and to industrial robotic-system workflows.}
\end{cvtable}

\cvsection{Education}
\subsection{Postgraduate Studies}
\begin{cvtable}
	\cvitem{Sep. 2021 -- Sep. 2024}{\color{cvsectioncolor}M.Sc. in Computer Engineering -- Artificial Intelligence \& Robotics}{\textbf{University of Padova, Italy}}{\textbf{Thesis}: Multi-Sensor Robotic Navigation Using a Teacher-Student Framework.\\
	\textbf{Grade}: 96/110 (approximately 1.9/4).\\
	\chartlabel{Machine Learning} \chartlabel{Computer Vision} \chartlabel{Robotics} \chartlabel{3D Data}}

	\cvitem{Dec. 2023 -- Sep. 2024}{\color{cvsectioncolor}Erasmus+ Exchange -- Thesis and Research Internship}{\textbf{University of Bonn, Germany}}{Research at the Humanoid Robots Lab under Prof. Emanuele Menegatti, Prof. Maren Bennewitz, and Murad Dawood.\\
	\chartlabel{TD3} \chartlabel{LiDAR} \chartlabel{Vision} \chartlabel{ROS/Gazebo}}
\end{cvtable}

\subsection{Undergraduate Study}
\begin{cvtable}
	\cvitem{Sep. 2015 -- Jun. 2020}{\color{cvsectioncolor}\small B.Sc. in Computer Software Engineering}{\textbf{University of Birjand, Iran}}{Coursework in C++, data structures, algorithm design, robotics, and image processing. \textbf{Grade}: 14.93/20.}
\end{cvtable}

\subsection{Other Training}
\begin{cvtable}
	\cvitem{2019 -- 2021}{\color{cvsectioncolor}\small Selected AI and Machine-Vision Certificates}{\textbf{DeepLearning.AI / Stemmer Imaging}}{TensorFlow for AI, machine learning and deep learning (2021); machine-vision solutions with HALCON and MERLIC (2019).}
\end{cvtable}

\newpage
\makebacksidebar

\cvsection{Selected University Projects}
\begin{cvtable}
	\cvitem{Dec. 2023}{\small Real-Time Camera Pose and Tumor Detection for Endoscopy}{Computer Vision}{Combined camera-motion tracking with HSV filtering, Canny edge detection, Hough Circle Transform, background subtraction, and morphological operations for real-time pose estimation and potential-tumor detection.}
	\cvitem{Oct. 2023}{\small TIAGO Mobile Robot Navigation and Obstacle Detection}{Robotics}{Implemented ROS navigation from a user-defined start to a destination and developed an Action Client/Server system for detecting and reporting movable obstacles from laser-sensor data.}
	\cvitem{Jun. 2023}{\small Image Colorization with Conditional WGAN}{Deep Learning}{Built a PyTorch conditional WGAN with a U-Net architecture and skip connections, trained on more than 50,000 grayscale images, and used conditional inputs and image-processing techniques for controlled colorization.}
	\cvitem{May 2023}{\small TinyPointNet for 3D Local Descriptors}{3D Data Processing}{Designed a compact PointNet variant for efficient local feature extraction from point clouds and evaluated it on diverse 3D datasets and benchmark tasks.}
	\cvitem{Jun. 2022}{\small Human Hands Detection and Segmentation}{Computer Vision}{Trained a custom YOLOv5 model with PyTorch and implemented C++ segmentation using contours, thresholding, and morphological operations.}
\end{cvtable}

\cvsection{Honours and Awards}
\begin{cvtable}
	\cvitem{}{Erasmus+ Traineeship Scholarship}{University of Padova}{}
	\cvitem{}{Regione Veneto Scholarship (2.5 years)}{University of Padova}{}
\end{cvtable}

\cvsection{Conference Experience}
\begin{cvtable}
	\cvitem{Mar. 2023}{Speaker -- 15th International Petroleum Technology Conference}{Bangkok, Thailand}{Presented ``Geological CO2 Leak Monitoring Using Joint Seismic-CSEM Inversion Based on CRIM and White Equations'' on behalf of a peer and communicated the technical content to an international audience.}
\end{cvtable}

\section{Teaching Experience}
\begin{cvtable}
	\cvitem{}{Teaching Assistant -- Differential Equations}{University of Birjand}{}
	\cvitem{}{Teaching Assistant -- General Mathematics I}{University of Birjand}{}
\end{cvtable}

\cvsection{References}
\begin{cvtable}
	\cvitem{Ref. 1}{Prof. AmirHossein MajidiRad}{Purdue University Northwest}{amajidir@pnw.edu}
	\cvitem{Ref. 2}{Prof. Emanuele Menegatti}{University of Padova}{emanuele.menegatti@unipd.it}
	\cvitem{Ref. 3}{Prof. Maren Bennewitz}{University of Bonn}{maren@cs.uni-bonn.de}
	\cvitem{Ref. 4}{Prof. Alberto Pretto}{University of Padova}{alberto.pretto@dei.unipd.it}
	\cvitem{Ref. 5}{Prof. Nasim Nasrabadi}{University of Birjand}{nasimnasrabadi@birjand.ac.ir}
\end{cvtable}

\newpage
\makethirdsidebar
\sidebarwidth=.08\paperwidth
\newgeometry{right=1cm,left=2cm,bottom=0.1cm,top=1cm}

\cvsection{Publications}
\subsection{Journals}
\begin{itemize}
{\small
\item Amani, A. M., Amani, S., and MajidiRad, A. H. ``A Unified Conditional Policy for Multi-Robot Navigation via LiDAR-to-Vision Distillation.'' Accepted for publication in the Robotics, Mechatronics and Intelligent Machines section of \textit{Machines} (MDPI), 2026.}
\end{itemize}
\vskip20pt
\subsection{Conference Presentation}
\begin{itemize}
{\small
\item ``Geological CO2 Leak Monitoring Using Joint Seismic-CSEM Inversion Based on CRIM and White Equations.'' Presented at the 15th International Petroleum Technology Conference in Bangkok on behalf of a peer, 2023.}
\end{itemize}

%\cvsignature

\end{document}
